\subsection{Evidence units and study selection}
The protocol uses three linked identities. An occurrence is one result as exported or captured from a run; a report is a specific full publication, including a later journal version; and a study is the distinct underlying evaluation. This distinction prevents counting a paper found in several databases more than once or treating a conference paper and its journal extension as independent experiments. DOI and title matching propose duplicate links, but ambiguous versions require a recorded decision. For example, the journal REEF report is preferred for GA-S001 and its earlier conference report remains linked history \cite{reef}. A different title or DOI is not sufficient by itself to prove a separate evaluation.

Title and abstract decisions use the protocol's scope criteria, with uncertainty advanced rather than forced into an exclusion. Retrieval attempts are recorded separately; inability to obtain a full text cannot be presented as an assessment of its evidence. At full text the reviewer applies one principal reason: E1 no inference-serving evidence; E2 no local multi-tenant sharing; E3 mechanism outside scope; E4 training only; E5 cluster only without separable local data; E6 non-primary or non-peer-reviewed report; E7 insufficient detail; E8 duplicate or less complete report; E9 period or language; E10 a documented other reason. These are decision categories, not an estimate of the unresolved queue's fate. The working set admits mixed training/inference studies only for separable inference outcomes and labels any transferred inference claim as such.

\subsection{Quality assessment and extraction}
All 43 working studies have a Q1--Q10 appraisal row with item scores and rationales. Each item receives 1, 0.5, or 0, with high (8--10), moderate (6--7.5), and low (below 6) reporting bands. Q1 checks the tenant question; Q2 the mechanism; Q3 hardware and software; Q4 workload; Q5 comparator; Q6 metrics; Q7 measurement controls; Q8 support for a result; Q9 limitations; and Q10 interpretability. The item-level matrix and checklist are in Appendix B. The total is a guide to how fully an evaluation can be interpreted; a high-scoring paper can still be indirect for this project's one-GPU decision or omit adversarial tests. No study was made eligible because it had a high score.

The extraction form records identity and report version; source type and access; GPU/SKU/count, OS, driver, runtime, partition/quota; model size and inference task; tenants, co-tenants, batch, arrival and measurement window; comparator; statistic and unit; bounded result; caveat; and page, figure or table anchor. Source-derived observation and reviewer inference are kept apart. The published CSV has separate performance interference, resource/memory, security, and fault-containment fields. A blank chart coordinate or unspecified driver is marked untranscribed where appropriate, rather than supplied from product defaults. The Appendix C cards reproduce a single bounded outcome and caveat per preferred report; the full workbook remains the evidence record for individual results.

\subsection{Synthesis plan and sensitivity}
The protocol calls for a structured narrative synthesis because workloads, accelerators, control layers, latency percentiles and resource budgets differ. Within-study matched comparisons can support a bounded direction or numeric effect in that paper's stratum. Across studies we compare conditions and mechanisms qualitatively; we do not average heterogeneous ratios or equate a server's GPU count to local occupancy. Findings are classified as direct within-study evidence, indirect transfer to the assigned decision, contextual mechanism description, or insufficient evidence. RQ1 uses the control-boundary taxonomy; RQ2 the metric and comparator audit; RQ3 the hardware/workload strata; RQ4 adverse examples and competing objectives; RQ5 a conditional decision table. A final sensitivity analysis must revisit recommendations after excluding moderate-score or weakly anchored evidence and after resolving the current access queue. \red{The working set has not passed that final sensitivity gate.}

\subsection{Search provenance and deviations}
The IEEE and Springer runs on 11 October and the later ACM and ScienceDirect exports on 14 October are separate events; an earlier blocked pilot is not a successful search. The exact run expressions, fields, filters, date, display observations and source links are rendered in Appendix A and kept in the two search-run CSVs. Springer and ACM use a broad master expression, while IEEE and ScienceDirect use family-specific expressions. The IEEE public search covered all metadata rather than the planned field-qualified command form. ScienceDirect split the master to satisfy its connector limit and omitted DNN/LLM in its shared inference block; this changes recall. ACM's displayed hit total and parsed exports diverge. The ScienceDirect temporal URL indicates an open-access facet whose use at the time of export is not established. These issues cannot be repaired by renaming imported rows. \red{A documented re-export or a formal coverage-limitation decision remains necessary.}
